\section{Implementation, certificates, and integration contracts}
\label{sec:implementation}

The package includes a runnable copy of the updated \code{fast} recognizer
and a patch against the pinned baseline. The patch changes or adds thirteen
paths, including two new mathematical modules, their tests, source-provenance
handling, CLI integration, documentation, and local certificate fixtures.
The existing license is preserved. This section describes the public
behavior that a repository integration should retain.

\subsection{Structured arithmetic and actual diagram construction}

The arithmetic API accepts an integer and a sequence of finite continued
fractions:
\begin{lstlisting}[language=Python]
from fastunknot.rational import (
    montesinos_certificate, verify_montesinos_certificate,
)

source = [[0, -3], [0, 5], [0, 7]]
certificate = montesinos_certificate(0, source)
assert certificate["status"] == "KNOTTED"
assert certificate["determinant"] == 1
assert verify_montesinos_certificate(0, source, certificate)
\end{lstlisting}
This call performs no PD expansion. It is the compressed-input interface
to Theorem~\ref{thm:mon:complete}. In particular, native Python integers
with many thousands of binary digits are supported by the arithmetic path.
The bit model in Proposition~\ref{prop:mon:complexity} describes those
integers, rather than a unit-cost model for arbitrarily large numbers.

For integration into the ordinary pipeline, the diagram constructor expands
the same source into a real plane diagram:
\begin{lstlisting}[language=Python]
from fastunknot import Diagram, recognize

d = Diagram.from_rational(0, [[0, -3], [0, 5], [0, 7]])
assert d.crossings == 15
answer = recognize(d)
assert answer.status == "KNOTTED"
assert answer.method == "montesinos-orbifold-quotient"

# Alias with the mathematical family name:
same = Diagram.from_montesinos(0, [[0, -3], [0, 5], [0, 7]])
\end{lstlisting}
The compiler implements horizontal twist composition, quarter rotation,
crossing reversal, and numerator closure. Rotation changes a slope to
\(-1/x\); reciprocal construction adds crossing reversal to obtain \(1/x\).
A lazy crossing-switch flag prevents repeated physical mirroring of all
already constructed rows at every reciprocal step. Its final output still
undergoes the ordinary spherical-embedding and component checks.

The corresponding JSON representation is:
\begin{lstlisting}
{
  "montesinos": {
    "e": 0,
    "tangles": [[0, -3], [0, 5], [0, 7]]
  }
}
\end{lstlisting}
The CLI accepts this through the ordinary input loader. The switch
\code{--no-rational} disables the new early decision, which is useful for
paired experiments on the very same expanded diagram. This switch does not
change what knot the supplied source represents.

The arithmetic stage recognizes both signs of the determinant-one
two-bridge condition. For instance,
\code{montesinos_certificate(0, [[1], [-1, -2]])}
describes \(N(T(1)-T(3/2))\) and returns \unk.
For \(P(-3,5,7)\), in contrast, the three exceptional denominators establish
nontriviality even though the determinant is one. The latter decision is
not an inference from the Alexander polynomial.

\subsection{Provenance is part of the correctness argument}

An arbitrary PD accompanied by an unrelated list of fractions is not
covered by the structured theorem. The implementation therefore attaches
immutable, detached rational source data only when its own checked
constructor creates the PD. The ordinary \code{Diagram.from_pd} constructor
does not invent or infer this provenance. Mutating the caller's original
list after construction does not mutate the stored source.

Mirroring preserves the connection between data and topology by changing
the source consistently. Pickle roundtrips replay construction and verify
the resulting source/PD relation. A guard also prevents a second
\code{Diagram.__init__} call from replacing the PD of an already initialized
object while leaving old provenance attached. These are ordinary supported
API invariants; the mathematics does not assume an attacker can rewrite
arbitrary internal interpreter objects.

This contract allows \code{recognize} to use a source without solving a
second, potentially difficult tangle-recognition problem. When the source
is absent, the existing PD pipeline is available. Automatic extraction of
rational structure from general diagrams is an explicit next research task.

\subsection{Input rejection and resource behavior}

The following cases are handled before a topological verdict is returned.
Every coefficient must be a literal integer; Python booleans are rejected
despite their being integer subclasses. Continued-fraction blocks must be
nonempty. Intermediate projective infinity is allowed, but a final
denominator zero is outside the supported source grammar. A final
multi-component numerator closure is outside the knot domain and is
rejected as input, rather than labeled a nontrivial knot.

The PD constructor defaults to a literal expansion ceiling of 100,000
crossings, checked before allocating the diagram. This policy protects the
expanded interface from exponentially large \(C\). It is separate from the
compressed arithmetic API, which does not need to allocate \(C\) crossings.
The current CLI uses the expanded constructor; the very large compressed
benchmarks below use the direct arithmetic API.

The existing recognition timeout and complex-object ceilings retain their
original meaning. Exhausting those limits yields \unknown, never \unk.
The separate expansion ceiling instead rejects construction with a
\code{DiagramError}; the CLI reports invalid input. Construction precedes
the recognition timer and is not covered by that timer. The exact arithmetic
branch is still subject to ordinary input processing and the chosen
pipeline budget; it does not create a universal real-time guarantee.

\subsection{Local occurrence certificates in arbitrary PDs}

The second new module is \code{fastunknot.tangle_obstruction}. Its direct
verification interface accepts a host diagram and a source-plus-map
certificate:
\begin{lstlisting}[language=Python]
from fastunknot.tangle_obstruction import (
    find_subtangle_obstruction,
    verify_subtangle_certificate,
    recognize_with_subtangles,
)

certificate, evidence, stats = find_subtangle_obstruction(host)
if certificate is not None:
    checked = verify_subtangle_certificate(host, certificate)
    assert checked["status"] == "KNOTTED"

# Optional search, followed by the existing recognizer if no witness is found.
answer = recognize_with_subtangles(host, seconds=3.0)
\end{lstlisting}
Here \code{host} is a normal validated \code{Diagram}; it need not carry
rational provenance. The example fixture files in \code{fast/examples}
are plain PDs, and their local certificates are separate files in
\code{fast/certificates}.

Each certificate contains exactly a version, a pattern source, a list of
crossing images, and a list of cyclic offsets:
\begin{lstlisting}
{
  "version": 1,
  "pattern": {"e": 0, "tangles": [[0, 3], [0, 3]]},
  "crossings": [...],
  "rotations": [0, 2, ...]
}
\end{lstlisting}
The ellipses above indicate the format; the delivered certificates contain
complete integer lists. Verification compiles the source or retrieves a
previously checked immutable template from the compilation cache. It then
checks the host and occurrence afresh: the crossing map, the cyclic order,
the overpassing pairs, every internal edge, the four distinct outgoing
edges, the common port face, and the two open strands. The checked template supplies the exact fractions
and obstruction. A stated knot verdict
in a JSON object is never itself a certificate.

The initial compiler accepts only \(e=0\), at least two finite nonintegral
rational summands, a connected crossing shadow, two open strands with no
closed component, and four distinct external cut edges at an occurrence.
The full arithmetic classifier has a broader source grammar; these two
interfaces should not be conflated.

The default catalogue consists of six templates: the sums with slopes
\((1/3,1/3)\), \((1/2,1/5)\), and \((1/3,1/4)\), and their mirrors.
Their literal sizes are six or seven crossings. Their respective
uncancelled congruence numerators are \(6,7,7\), with products \(9,10,12\).
Each fails the \(\pm1\) congruence and has at most one even denominator.

Connectedness makes matching unusually simple. Once a root crossing image
and an allowed half-turn are selected, every internal edge forces its
neighbor's image and rotation. The search does not enumerate all
\(n^m\) crossing maps. For an already compiled \(m\)-crossing pattern it
uses the \(O(nm)\) forced-propagation algorithm proved above.
The public call also sorts labels during validation, adding
\(O(n\log(n+2))\) word operations, plus source reading and compilation.
These deterministic word bounds can be realized with indexed arrays.
The Python implementation uses dictionaries and sets in some checks;
its operation accounting uses the usual constant-time table model.
No worst-case bit bound for adversarial interpreter hash-table behavior
is being inferred from a wall-clock measurement.

The wrapper is opt-in. An inconclusive search continues to the existing
pipeline, and a failed time budget has the same \unknown\ interpretation.
Successful local rejection is independent of the exterior's knot
polynomial or homology. An unknottable local pattern, however, supplies
no positive verdict about its actual exterior.

\subsection{What the integration patch does and does not contain}

The patch includes \code{rational.py}, \code{tangle_obstruction.py},
the new tests and fixtures, modifications to the diagram class and
recognition entry points, and updated documentation. It can be reviewed
and applied from the ProveIt repository root; the exact commands and
baseline checks are in \code{integration/README.md}.

The archive also contains the entire runnable \code{fast} directory,
the article sources and PDF, example-source generators, independent
diagnostics, raw measurements, and manifests. The theoretical exhaustive
enumerator of Proposition~\ref{prop:mon:bounded-discovery} is not silently
presented as shipping software. The implementation supplies the required
fixed-pattern search primitive; a practical enumerator should be a separate
reviewed contribution with source deduplication and explicit budgets.
